\chapter{量子アルゴリズムへの接続}
\begin{chaptergoal}
\cgitem 線形化の後にも残る演算子構成と読み出しの課題を知る。\\
\cgitem ブロックエンコーディングの規格化係数を理解する。\\
\cgitem 漸近的な高速化の主張を、その前提とともに読む。
\end{chaptergoal}

\section{巨大な線形系を量子状態へ}
Carleman系は線形だが、古典的に全成分を保存すれば $D_{N_C}$ の増大を避けられない。
量子計算では長いベクトルを振幅に符号化できるため、状態の添字自体には対数個の量子ビットを使える。
ただし入力準備、演算子の実装、時間発展、観測量の読み出しも効率的でなければ、状態表現の小ささは計算量の優位性につながらない。
本章は提供された論文解説に基づく理論の紹介であり、現行コードの機能ではない。

\section{ブロックエンコーディングの役割}
非ユニタリ行列 $M$ を直接量子ゲートとすることはできない。
補助量子ビットを加えたユニタリ $U_M$ の一部として
\begin{equation}
(\bra{0}^{\otimes a}\otimes I)U_M(\ket{0}^{\otimes a}\otimes I)
\simeq M/\alpha_M
\end{equation}
を実現する方法がブロックエンコーディングである。
$\alpha_M$ は少なくとも $\normtwo{M}$ 程度が必要で、アルゴリズムの反復回数などに影響する。
行列が大きいことだけでなく、どの構造を使って小さな規格化係数で構成できるかが重要になる。

遮蔽された電場の実空間相互作用は距離とともに減衰する。
階層的ブロックエンコーディングは、近距離を疎な相互作用として、遠距離を大きなブロックとして扱い、距離ごとの強度を構成へ反映する。
複数のユニタリの線形結合を実装するLCU法を使い、各階層を合成する。
提供資料では原論文 Eq.~(150) の結果として
\begin{equation}
\alpha_{\widetilde F_2}=O(\sqrt{N_FN_H}),\qquad
\alpha_{\widetilde F_1}=O(N_F\sqrt{N_H})
\end{equation}
が紹介される。固定した物理パラメータなどの依存性は省略されている。
単に現在の密な $F_2$ 配列を量子メモリへ格納すれば得られる評価ではない。

\section{履歴状態と観測量の推定}
量子ODEソルバは、時刻レジスタを含む履歴状態を準備する。
概念的には
\begin{equation}
\ket{\chi}\propto\sum_{s=0}^{M-1}
\normtwo{\bm z(s\Delta t)}\ket{\bm z(s\Delta t)}\ket{s},
\qquad \ket{\bm z(t)}=\frac{1}{\normtwo{\bm z(t)}}\sum_i z_i(t)\ket{i}
\end{equation}
と書ける。状態ベクトルの全成分を測定して復元すると、出力数に応じた大きな費用がかかる。
論文での目標は、時空間平均された非平衡運動エネルギーモーメントという指定された量の推定である。
高次積分近似と振幅推定を組み合わせ、時間平均と読み出しにかかる精度依存性も考慮する。

\begin{notebox}
現行コードが図示する $\mathcal F=\uu^\dagger P\uu/2$ は安定性評価に用いる二次形式である。
量子アルゴリズムの出力として紹介される運動エネルギーモーメントと、名称やグラフを取り違えない。
\end{notebox}

\section{計算量の比較を読む}
提供資料は、3D+3D系の代表的な比較として次の演算量を挙げている。
\begin{align}
\mathcal C_{\mathrm{op}}&=\widetilde O\left(
T N_F^4N_H^{7/2}\epsilon^{-1/p}\right),\\
\mathcal Q_{\mathrm{op}}&=\widetilde O\left(
N_FN_H^{1/2}\operatorname{polylog}(T/\epsilon)\epsilon^{-1}\right).
\end{align}
$\widetilde O$ は対数因子を省略し、$p$ は比較対象の古典的な陽的時間積分法の次数を表す。
古典メモリは $O(N_F^3N_H^3)$、量子ビット数は
\[
\widetilde O\left(\log(N_FN_H^{1/2}T)\log(1/\epsilon)\right)
\]
と紹介される。所定の弱非線形条件、初期状態準備、演算子アクセス、観測量の仮定に基づく漸近評価である。
1D1V実装の計算量を直接表す式ではない。

空間が高い正則性を持ち、速度方向のHermite正則性 $s_H=1$、$p=4$ の例では
\[
\mathcal Q_{\mathrm{op}}=\widetilde O(\epsilon^{-3/2}),\quad
\mathcal C_{\mathrm{op}}=\widetilde O(\epsilon^{-15/4}),\quad
\mathcal C_{\mathrm{op}}=\widetilde O(\mathcal Q_{\mathrm{op}}^{5/2})
\]
となる。これが「二次を超える高速化」という比較の意味である。
実際の定数、誤り訂正、入力データの準備費、比較する古典手法を含む実行時間の優位性は、この指数比較だけでは決まらない。

\section{保証される物理領域の狭さ}
提供資料には、原論文の代表条件 $\nuu=10^{-2}$ とDebyeスケールの分解能の下で、
無次元電位 $\varphi_{\max}$ の保証上限が単一Fourierモードで $7.12\times10^{-8}$、Gaussian密度摂動で $2.08\times10^{-5}$ と記載される。
同資料のトカマク炉心の典型値 $10^{-3}$〜$10^{-2}$ との隔たりは大きい。
これらは本プロジェクトの $\nuu=0.2$ の結果から得た閾値ではなく、一般の破綻点でもない。
また $\varphi$ と本書の $\phi$ の物理単位・温度基準を確認せずに直接比較してはいけない。
保証条件の保守性、磁場や現実的衝突演算子への拡張、より強い非線形現象への対応が研究課題として残る。
